% Figure: ATP yield per gram of substrate for the main cellular fuels.
\begin{figure}[htbp]
\centering
\begin{tikzpicture}
\begin{axis}[
  width=10.5cm, height=6cm,
  ybar,
  bar width=14pt,
  ymin=0, ymax=0.50,
  ylabel={ATP yield (mol/g substrate)},
  symbolic x coords={Glucose (anaerobic), Glucose (aerobic), Palmitate (aerobic), Alanine (aerobic), Ethanol (aerobic)},
  xtick=data,
  xticklabel style={font=\scriptsize, rotate=22, anchor=east},
  ylabel style={font=\small},
  ytick={0,0.1,0.2,0.3,0.4,0.5},
  ymajorgrids,
  major grid style={draw=engblack!15},
  nodes near coords,
  nodes near coords align={vertical},
  every node near coord/.append style={font=\scriptsize},
]
% Yields per gram, derived from data/atp_yield_per_substrate.csv
% mass yields (mol ATP / g):
%   glucose (anaerobic, fermentation): 2 ATP / 180 g/mol = 0.0111
%   glucose (aerobic): 32 / 180 = 0.178
%   palmitate (C16): 106 / 256 = 0.414
%   alanine: ~13 / 89 = 0.146
%   ethanol (aerobic): ~14 / 46 = 0.304
\addplot[draw=engblue, fill=engblue!30]
  coordinates {
    (Glucose (anaerobic),0.011)
    (Glucose (aerobic),0.178)
    (Palmitate (aerobic),0.414)
    (Alanine (aerobic),0.146)
    (Ethanol (aerobic),0.304)};
\end{axis}
\end{tikzpicture}
\caption[ATP yield per gram of fuel]{Per-gram molar ATP yield for
common cellular substrates. Anaerobic glycolysis to lactate from
glucose nets 2 ATP per glucose; aerobic complete oxidation nets $\sim$
32. Fatty acids (palmitate, C$_{16}$) yield more ATP per gram (longer
hydrocarbon chain, more electrons donated to the chain). Per-gram
ethanol yield is high but small in absolute terms because ethanol is
a small molecule. Per-gram fat is the densest cellular fuel by a
factor of $\sim 2$ over carbohydrate, matching the dietary value of
$9.4$ vs $4$ kcal/g. Values from \texttt{data/atp\_yield\_per\_substrate.csv}.}
\label{fig:vol06:ch02:fuel-yield}
\end{figure}
